\section{Referenzierte Schlussformen f\"ur Modellargumente}

In diesem Abschnitt sind die Aussagen \(M\models\varphi\) und
\(T\vdash\varphi\) Aussagen der zuvor festgelegten Metatheorie.
Auch auf dieser Ebene werden die Schlussregeln aus Band~1 angewendet.
Die letzte Spalte nennt jeweils die Regel oder den zuvor bewiesenen
Satz und die benutzten Zeilen. Eigenvariablen werden in jedem
Subbeweis frisch gew\"ahlt.

\FormulaDefDeltaK[Ausgeschriebene Modell- und Erweiterungsnotation]
{\begin{gathered}
 T+\varphi\coloneqq T\cup\{\varphi\},\\
 M\models T\ \coloneqq\
 \forall\theta\,(\theta\in T\longrightarrow M\models\theta),\\
 T\nvdash\varphi\ \coloneqq\ \neg(T\vdash\varphi)
 \end{gathered}}
{B48ModelReasonNotation}{
 \DeltaRow{Theorie}{T\subseteq\operatorname{Sent}(\mathcal L)}
 \DeltaRow{Satz}{\varphi\in\operatorname{Sent}(\mathcal L)}
 \DeltaRow{Struktur}{M}
}
Diese Definition schreibt die bereits benutzten Abk\"urzungen aus;
sie f\"uhrt kein zus\"atzliches Mengenaxiom ein. Insbesondere wird
die Erf\"ullung einer erweiterten Theorie nicht als neue
Schlussregel vorausgesetzt, sondern im folgenden Satz hergeleitet.

\FormulaThmDeltaK[Mitgliedschaft in einer Theorieerweiterung]
{\theta\in T+\varphi\ \longleftrightarrow\
 (\theta\in T\lor\theta=\varphi)}
{B48TheoryExtensionMember}{
 \DeltaRow{Theorie}{T\subseteq\operatorname{Sent}(\mathcal L)}
 \DeltaRow{S\"atze}{\theta,\varphi\in\operatorname{Sent}(\mathcal L)}
}
\begin{tabproof}
 \proofstep{1}{\theta\in T+\varphi}{\rA}
 \proofstep{1}{\theta\in T\cup\{\varphi\}}
 {\FormulaRefAuto{B48ModelReasonNotation}{1}}
 \proofstep{1}{\theta\in T\lor\theta\in\{\varphi\}}
 {\FormulaRefAuto{x\in A\cup B\eqvdash x\in A\lor x\in B}{2}}
 \proofstep{4}{\theta\in T}{\rA}
 \proofstep{4}{\theta\in T\lor\theta=\varphi}
 {\RuleRef{OI1}{\lor I_1}\,4}
 \proofstep{6}{\theta\in\{\varphi\}}{\rA}
 \proofstep{6}{\theta=\varphi}
 {\FormulaRefAuto{x\in\{a\}\eqvdash x=a}{6}}
 \proofstep{6}{\theta\in T\lor\theta=\varphi}
 {\RuleRef{OI2}{\lor I_2}\,7}
 \proofstep{1}{\theta\in T\lor\theta=\varphi}
 {\RuleRef{OE}{\lor E}\,3,4\!:\!5,6\!:\!8}
 \proofstep{}{\theta\in T+\varphi\rightarrow
   (\theta\in T\lor\theta=\varphi)}
 {\RuleRef{RI}{\rightarrow I}\,1\!:\!9}
 \proofstep{11}{\theta\in T\lor\theta=\varphi}{\rA}
 \proofstep{12}{\theta\in T}{\rA}
 \proofstep{12}{\theta\in T\lor\theta\in\{\varphi\}}
 {\RuleRef{OI1}{\lor I_1}\,12}
 \proofstep{14}{\theta=\varphi}{\rA}
 \proofstep{14}{\theta\in\{\varphi\}}
 {\FormulaRefAuto{x\in\{a\}\eqvdash x=a}{14}}
 \proofstep{14}{\theta\in T\lor\theta\in\{\varphi\}}
 {\RuleRef{OI2}{\lor I_2}\,15}
 \proofstep{11}{\theta\in T\lor\theta\in\{\varphi\}}
 {\RuleRef{OE}{\lor E}\,11,12\!:\!13,14\!:\!16}
 \proofstep{11}{\theta\in T\cup\{\varphi\}}
 {\FormulaRefAuto{x\in A\cup B\eqvdash x\in A\lor x\in B}{17}}
 \proofstep{11}{\theta\in T+\varphi}
 {\FormulaRefAuto{B48ModelReasonNotation}{18}}
 \proofstep{}{(\theta\in T\lor\theta=\varphi)
   \rightarrow\theta\in T+\varphi}
 {\RuleRef{RI}{\rightarrow I}\,11\!:\!19}
 \proofstep{}{\theta\in T+\varphi\leftrightarrow
   (\theta\in T\lor\theta=\varphi)}
 {\RuleRef{LRI}{\leftrightarrow I}\,10,20}
\end{tabproof}

\FormulaThmDeltaK[Modell einer Theorieerweiterung]
{\begin{gathered}
 M\models T+\varphi\ \longleftrightarrow\\
 (M\models T)\land(M\models\varphi)
 \end{gathered}}
{B48ModelTheoryExtension}{
 \DeltaRow{Theorie}{T\subseteq\operatorname{Sent}(\mathcal L)}
 \DeltaRow{Satz}{\varphi\in\operatorname{Sent}(\mathcal L)}
 \DeltaRow{Struktur}{M}
}
\begin{tabproof}
 \proofstep{1}{M\models T+\varphi}{\rA}
 \proofstep{1}{\forall\theta\,(\theta\in T+\varphi
   \rightarrow M\models\theta)}
 {\FormulaRefAuto{B48ModelReasonNotation}{1}}
 \proofstep{3}{\theta\in T}{\rA}
 \proofstep{3}{\theta\in T\lor\theta=\varphi}
 {\RuleRef{OI1}{\lor I_1}\,3}
 \proofstep{3}{\theta\in T+\varphi}
 {\FormulaRefAuto{B48TheoryExtensionMember}{4}}
 \proofstep{1}{\theta\in T+\varphi\rightarrow M\models\theta}
 {\RuleRef{UE}{\forall E}\,2}
 \proofstep{1,3}{M\models\theta}
 {\RuleRef{RE}{\rightarrow E}\,6,5}
 \proofstep{1}{\theta\in T\rightarrow M\models\theta}
 {\RuleRef{RI}{\rightarrow I}\,3\!:\!7}
 \proofstep{1}{\forall\theta\,(\theta\in T\rightarrow M\models\theta)}
 {\RuleRef{UI}{\forall I}\,8}
 \proofstep{1}{M\models T}
 {\FormulaRefAuto{B48ModelReasonNotation}{9}}
 \proofstep{}{\varphi=\varphi}{\RuleRef{II}{=I}}
 \proofstep{}{\varphi\in T\lor\varphi=\varphi}
 {\RuleRef{OI2}{\lor I_2}\,11}
 \proofstep{}{\varphi\in T+\varphi}
 {\FormulaRefAuto{B48TheoryExtensionMember}{12}}
 \proofstep{1}{\varphi\in T+\varphi\rightarrow M\models\varphi}
 {\RuleRef{UE}{\forall E}\,2}
 \proofstep{1}{M\models\varphi}
 {\RuleRef{RE}{\rightarrow E}\,14,13}
 \proofstep{1}{(M\models T)\land(M\models\varphi)}
 {\RuleRef{AI}{\land I}\,10,15}
 \proofstep{}{\begin{gathered}M\models T+\varphi\rightarrow\\
   ((M\models T)\land(M\models\varphi))\end{gathered}}
 {\RuleRef{RI}{\rightarrow I}\,1\!:\!16}
 \proofstep{18}{(M\models T)\land(M\models\varphi)}{\rA}
 \proofstep{18}{M\models T}{\RuleRef{AE1}{\land E_1}\,18}
 \proofstep{18}{M\models\varphi}{\RuleRef{AE2}{\land E_2}\,18}
 \proofstep{18}{\forall\theta\,(\theta\in T\rightarrow M\models\theta)}
 {\FormulaRefAuto{B48ModelReasonNotation}{19}}
 \proofstep{22}{\theta\in T+\varphi}{\rA}
 \proofstep{22}{\theta\in T\lor\theta=\varphi}
 {\FormulaRefAuto{B48TheoryExtensionMember}{22}}
 \proofstep{24}{\theta\in T}{\rA}
 \proofstep{18}{\theta\in T\rightarrow M\models\theta}
 {\RuleRef{UE}{\forall E}\,21}
 \proofstep{18,24}{M\models\theta}
 {\RuleRef{RE}{\rightarrow E}\,25,24}
 \proofstep{27}{\theta=\varphi}{\rA}
 \proofstep{}{\theta=\theta}{\RuleRef{II}{=I}}
 \proofstep{27}{\varphi=\theta}{\RuleRef{IE}{=E}\,27,28}
 \proofstep{18,27}{M\models\theta}{\RuleRef{IE}{=E}\,29,20}
 \proofstep{18,22}{M\models\theta}
 {\RuleRef{OE}{\lor E}\,23,24\!:\!26,27\!:\!30}
 \proofstep{18}{\theta\in T+\varphi\rightarrow M\models\theta}
 {\RuleRef{RI}{\rightarrow I}\,22\!:\!31}
 \proofstep{18}{\forall\theta\,(\theta\in T+\varphi
   \rightarrow M\models\theta)}
 {\RuleRef{UI}{\forall I}\,32}
 \proofstep{18}{M\models T+\varphi}
 {\FormulaRefAuto{B48ModelReasonNotation}{33}}
 \proofstep{}{\begin{gathered}((M\models T)\land(M\models\varphi))\\
   \rightarrow M\models T+\varphi\end{gathered}}
 {\RuleRef{RI}{\rightarrow I}\,18\!:\!34}
 \proofstep{}{\begin{gathered}M\models T+\varphi\leftrightarrow\\
   ((M\models T)\land(M\models\varphi))\end{gathered}}
 {\RuleRef{LRI}{\leftrightarrow I}\,17,35}
\end{tabproof}

\FormulaThmDeltaK[Ein Modell verhindert eine Widerlegung]
{M\models T,\ M\models\varphi\ \vdash\ T\nvdash\neg\varphi}
{B48ModelBlocksRefutation}{
 \DeltaRow{Theorie}{T\subseteq\operatorname{Sent}(\mathcal L)}
 \DeltaRow{Satz}{\varphi\in\operatorname{Sent}(\mathcal L)}
 \DeltaRow{Struktur}{M}
}
\begin{tabproof}
 \proofstep{1}{M\models T}{\rA}
 \proofstep{2}{M\models\varphi}{\rA}
 \proofstep{3}{T\vdash\neg\varphi}{\rA}
 \proofstep{1,3}{M\models\neg\varphi}
  {\FormulaRefAuto{B48Soundness}{3,1}}
 \proofstep{1,2,3}{\bot}
  {\FormulaRefAuto{B48SentenceContradiction}{2,4}}
 \proofstep{1,2}{\neg(T\vdash\neg\varphi)}{\rCI{3,5}}
 \proofstep{1,2}{T\nvdash\neg\varphi}
  {\FormulaRefAuto{B48ModelReasonNotation}{6}}
\end{tabproof}

\FormulaThmDeltaK[Ein Gegenmodell verhindert einen Beweis]
{M\models T,\ M\models\neg\varphi\ \vdash\ T\nvdash\varphi}
{B48CountermodelBlocksProof}{
 \DeltaRow{Theorie}{T\subseteq\operatorname{Sent}(\mathcal L)}
 \DeltaRow{Satz}{\varphi\in\operatorname{Sent}(\mathcal L)}
 \DeltaRow{Struktur}{M}
}
\begin{tabproof}
 \proofstep{1}{M\models T}{\rA}
 \proofstep{2}{M\models\neg\varphi}{\rA}
 \proofstep{3}{T\vdash\varphi}{\rA}
 \proofstep{1,3}{M\models\varphi}
  {\FormulaRefAuto{B48Soundness}{3,1}}
 \proofstep{1,2,3}{\bot}
  {\FormulaRefAuto{B48SentenceContradiction}{4,2}}
 \proofstep{1,2}{\neg(T\vdash\varphi)}{\rCI{3,5}}
 \proofstep{1,2}{T\nvdash\varphi}
  {\FormulaRefAuto{B48ModelReasonNotation}{6}}
\end{tabproof}

\FormulaThmDeltaK[Relative Konsistenz verhindert Widerlegbarkeit]
{\BXLVIIICon(T+\varphi)\longrightarrow T\nvdash\neg\varphi}
{B48RelativeNonRefutation}{
 \DeltaRow{Sprache}{\mathcal L\text{ abz\"ahlbar, relational}}
 \DeltaRow{Theorie}{T\subseteq\operatorname{Sent}(\mathcal L)}
 \DeltaRow{Satz}{\varphi\in\operatorname{Sent}(\mathcal L)}
}
\begin{tabproof}
 \proofstep{1}{\BXLVIIICon(T+\varphi)}{\rA}
 \proofstep{1}{\exists M\,(M\models T+\varphi\land |M|\le\aleph_0)}
 {\FormulaRefAuto{B48Completeness}{1}}
 \proofstep{3}{M\models T+\varphi\land |M|\le\aleph_0}{\rA}
 \proofstep{3}{M\models T+\varphi}{\RuleRef{AE1}{\land E_1}\,3}
 \proofstep{3}{(M\models T)\land(M\models\varphi)}
 {\FormulaRefAuto{B48ModelTheoryExtension}{4}}
 \proofstep{3}{M\models T}{\RuleRef{AE1}{\land E_1}\,5}
 \proofstep{3}{M\models\varphi}{\RuleRef{AE2}{\land E_2}\,5}
 \proofstep{3}{T\nvdash\neg\varphi}
 {\FormulaRefAuto{B48ModelBlocksRefutation}{6,7}}
 \proofstep{1}{T\nvdash\neg\varphi}
 {\RuleRef{EE}{\exists E}\,2,3\!:\!8}
 \proofstep{}{\BXLVIIICon(T+\varphi)\rightarrow T\nvdash\neg\varphi}
 {\RuleRef{RI}{\rightarrow I}\,1\!:\!9}
\end{tabproof}

\FormulaThmDeltaK[Relative Konsistenz der Negation verhindert Beweisbarkeit]
{\BXLVIIICon(T+\neg\varphi)\longrightarrow T\nvdash\varphi}
{B48RelativeNonProof}{
 \DeltaRow{Sprache}{\mathcal L\text{ abz\"ahlbar, relational}}
 \DeltaRow{Theorie}{T\subseteq\operatorname{Sent}(\mathcal L)}
 \DeltaRow{Satz}{\varphi\in\operatorname{Sent}(\mathcal L)}
}
\begin{tabproof}
 \proofstep{1}{\BXLVIIICon(T+\neg\varphi)}{\rA}
 \proofstep{1}{\exists M\,(M\models T+\neg\varphi\land |M|\le\aleph_0)}
 {\FormulaRefAuto{B48Completeness}{1}}
 \proofstep{3}{M\models T+\neg\varphi\land |M|\le\aleph_0}{\rA}
 \proofstep{3}{M\models T+\neg\varphi}{\RuleRef{AE1}{\land E_1}\,3}
 \proofstep{3}{(M\models T)\land(M\models\neg\varphi)}
 {\FormulaRefAuto{B48ModelTheoryExtension}{4}}
 \proofstep{3}{M\models T}{\RuleRef{AE1}{\land E_1}\,5}
 \proofstep{3}{M\models\neg\varphi}{\RuleRef{AE2}{\land E_2}\,5}
 \proofstep{3}{T\nvdash\varphi}
 {\FormulaRefAuto{B48CountermodelBlocksProof}{6,7}}
 \proofstep{1}{T\nvdash\varphi}
 {\RuleRef{EE}{\exists E}\,2,3\!:\!8}
 \proofstep{}{\BXLVIIICon(T+\neg\varphi)\rightarrow T\nvdash\varphi}
 {\RuleRef{RI}{\rightarrow I}\,1\!:\!9}
\end{tabproof}
